\section{Conclusion}

This project ported Granite-4.0-h-tiny and Granite-4.0-h-small to Tenstorrent Wormhole hardware and benchmarked them against an NVIDIA A100 and CPU. The tiny model (4 chips) achieves 10.65--10.69 tok/s, which is ${\approx}18\%$ faster than the A100 (8.5--9.0 tok/s) and ${\approx}4\times$ faster than CPU. The small model (8 chips) achieves 5.84--5.89 tok/s, which is ${\approx}20\%$ below the A100 (7.1--8.2 tok/s) because reading the larger model's weights across 8 chips takes more time than the A100 can match at this scale, but is still ${\approx}9\times$ faster than CPU.

The key enabler for the tiny model result is decode trace: recording the full set of operations once and replaying it for every token removes the per-token overhead of sending commands from the host to the device. A critical implementation detail is that trace recording must happen outside the generation loop — recording inside the loop causes all state updates to lag by one step, producing incorrect output. This was fixed by capturing the trace in a separate setup step before generation begins.

Two custom fused kernels (\texttt{ssm\_update} and \texttt{conv1d\_decode}) reduced the number of memory reads per Mamba2 step. At batch size 1 the gains were modest (${\approx}1$--2\%) because most time is spent reading weights from memory rather than doing arithmetic, but the gains are expected to be larger at higher batch sizes. Prefill chunk size tuning showed that $C{=}256$ gives the best decode throughput.

The main challenges encountered during porting — including weight quantisation, SSM state management across chunked prefill, expert-parallel weight sharding, and collective operation ordering — are documented in Section~\ref{sec:challenges} as a reference for future ports of similar hybrid models to Wormhole hardware.

Several directions remain open for future work. All benchmarks in this project generate one token at a time (batch size~1), which is a particularly demanding case for throughput because the hardware spends most of its time waiting for weight data to arrive from memory rather than doing computation. At higher batch sizes, multiple requests are processed together, so the same weight data serves several tokens at once and the memory bottleneck eases. Measuring how throughput scales with batch size would show where Wormhole hardware becomes most competitive, and would also be where the custom fused kernels are expected to give larger gains~\citep{dao2024transformers}. Power consumption is another dimension not explored here. The Galaxy 32 system draws 6--8\,kW on average~\citep{tenstorrent_galaxy} across all 32 chips, and the A100 PCIe card is rated at 300\,W~\citep{nvidia2020a100}. The two systems are not directly comparable on power since the tiny model uses only 4 of the 32 chips and the small model uses 8. Measuring actual wall power under load for a fair per-chip comparison is left for future work.

A second line of future work concerns inference speed at batch size~1 specifically. Speculative decoding~\citep{leviathan2023fast} is a technique that uses two models together: a small, fast model generates several candidate tokens quickly, and the main model then checks all of them in one step rather than generating tokens one at a time. If the main model agrees with the candidates, multiple tokens are accepted from that single step, increasing the number of tokens produced per second. Applying this to the Wormhole port could improve single-request throughput beyond what the trace optimisation alone achieves. Finally, the solutions developed in this project --- for SSM state management, expert-parallel sharding, and multi-chip communication ordering --- were designed around the specific structure of Granite. Porting a second model that combines recurrent layers with sparse expert routing would test whether these solutions generalise or whether further adaptation is needed.
